\documentclass[10pt,a4paper]{amsart}
\usepackage[margin=19mm]{geometry}
\usepackage{amsmath,amssymb,mathtools}
\usepackage[scr=rsfs,cal=euler]{mathalfa}
\usepackage{xfrac}
\usepackage[unicode,hidelinks,bookmarks=false]{hyperref}
\newcommand{\ie}{i.e.\ }
\newcommand{\cf}{cf.\ }
\newcommand{\resp}{respectively\ }
\newcommand{\Subsection}{\subsection}
\providecommand{\nobreakdash}{\nobreak\hbox{-}}
\setlength{\emergencystretch}{2em}
\multlinegap0pt
\usepackage{bm}
\usepackage{bbm}
\usepackage{dsfont}
\usepackage{xspace}
\usepackage{cancel}
\usepackage{mathtools}
\usepackage{amsthm}
\makeatletter
\@ifundefined{defin}{\newtheorem{defin}{Defin}}{}
\@ifundefined{corollary}{\newtheorem{corollary}[defin]{Corollary}}{}
\@ifundefined{lemma}{\newtheorem{lemma}[defin]{Lemma}}{}
\makeatother
\makeatletter
\newcommand{\C}{\mathbb{C}}
\newcommand{\N}{\mathbb{N}}
\expandafter\ifx\csname definition\endcsname\relax
\newenvironment{definition}{\begin{defin}\rm}{\end{defin}}
\fi
\expandafter\ifx\csname example\endcsname\relax
\newenvironment{example}{\begin{exa}\rm}{\end{exa}}
\fi
\expandafter\ifx\csname examples\endcsname\relax
\newenvironment{examples}{\begin{exas}\rm}{\end{exas}}
\fi
\expandafter\ifx\csname proof\endcsname\relax
\newenvironment{proof}
{\noindent{\it Proof.}}{\hfill $\Box$\par\vspace{2.5mm}}
\fi
\expandafter\ifx\csname remark\endcsname\relax
\newenvironment{remark}
%{\par\vspace{2.5mm}\noindent{\it Remark.}}{\par\vspace{2.5mm}}
\fi
\makeatother
\title[$Q$-difference analogue of the Stothers-Mason theorem]{$Q$-difference analogue of the Stothers-Mason theorem}
\author{Jian-Tang Lu, Xing-Xing Lu, Zhi-Tao Wen}
\date{Source: arXiv:2605.27876v1 (CC BY 4.0)}
\begin{document}
\maketitle

\noindent\emph{Source and licence.} $Q$-difference analogue of the Stothers-Mason theorem. Version arXiv:2605.27876v1 (CC BY 4.0), distributed under the Creative Commons Attribution 4.0 International licence, as recorded in the arXiv licence metadata for this exact version and shown on the abstract page https://arxiv.org/abs/2605.27876v1. Full rights notice: licenses/arxiv-2605.27876-CC-BY-4.0.txt.\par\smallskip

\noindent\textit{Editorial scope.} Counted unit: the Lemma whose source label is \texttt{q-prime} in the article (source0.tex lines 355--359, source label \texttt{q-prime}), delivered together with the article's own complete original proof of it (source0.tex lines 361--402, one \texttt{proof} environment of 42 lines). No mathematics is written, repaired or re-derived: statement and proof are the article's own text line for line. Required context retained uncounted: n-1.cor (lines 280--283). Every definition, display and label of those lines is retained unchanged. Omitted from this package and not redistributed: 1--279, 284--354, 360--360, 403--793. Nothing inside a retained excerpt is omitted or altered: every retained body is delivered exactly as the article's lines are written, so no omission receipt of any kind accompanies this package. Citations: the retained excerpts carry no citation command at all, so no bibliography entry of the article is transcribed and no reference is redistributed. Reference adaptation: none was needed and none was made; every reference inside the delivered unit resolves inside it, so no unresolved reference or citation is produced. Printed slips kept literal: no printed word, symbol, hypothesis or number of the counted statement or of its proof was altered. Layout and class: the article's own class is \texttt{article} with packages including inputenc, amsfonts, latexsym, amsmath, amssymb, multirow, eepic, graphicx, color, tikz, geometry; the wrapper is the lane's pinned \texttt{amsart} editorial wrapper at 10pt on A4, which restates the theorem-like environments the retained text needs and adds the article's own macro definitions that the retained text needs (\texttt{\textbackslash C}, \texttt{\textbackslash N}), with extra wrapper packages (bm, bbm, dsfont, xspace, cancel, mathtools). The delivered numbering is the unit's own. Assets: no retained span contains a figure environment, a graphics-inclusion command, a PDF-page inclusion, a table or a TikZ picture; no image file is copied into this package, the package declares no asset dependency and no image file is redistributed with it. Licence: version arXiv:2605.27876v1 of this article is distributed under the Creative Commons Attribution 4.0 International licence, as recorded in the arXiv licence metadata for this exact version and shown on the abstract page https://arxiv.org/abs/2605.27876v1. A full rights notice is delivered at licenses/arxiv-2605.27876-CC-BY-4.0.txt. Characters: every retained excerpt is pure ASCII, so the delivered engine, XeLaTeX with its default Latin Modern text font, reports no missing character.
\begin{corollary}\label{n-1.cor}
Let $f$ be analytic in the domain $G\subset\C$, and let $n\in\N^+$ and $q\in\C\setminus\{0\}$. Suppose that $z_0, qz_0,\ldots, q^nz_0\in G$.
If $z_0$ is an $a$-point of $f$ with $q$-weight $n$ in $G$, then $z_0$ is a zero of $D_q f$ with $q$-weight $n-1$.
\end{corollary}

\begin{lemma}\label{q-prime}
Suppose that $q\in\C\setminus\{0\}$, and $f$ and $g$ are relatively $q$-prime entire functions. Then ${\rm gcd}(f,D_q f)$
and ${\rm gcd}(g,D_q g)$ are also relatively $q$-prime. In addition, ${\rm gcd}(f,D_q f)$
and ${\rm gcd}(g,D_q g)$ are relatively prime.
\end{lemma}

\begin{proof}
Assume that ${\rm gcd}(f,D_q f)$ and ${\rm gcd}(g,D_q g)$ have a common $q$-divisor.
There exist $z_1,z_2\in\C$, $m,n\in\N$ and entire functions $F,G$ such that
    \begin{equation}\label{fDf.eq}
    {\rm gcd}(f,D_q f)=[z-z_1]_q^mF(z),
    \end{equation}
and
    \begin{equation}\label{gDg.eq}
    {\rm gcd}(g,D_q g)=[z-z_2]_q^nG(z).
    \end{equation}
Without loss of generality, we suppose that $z_2=q^mz_1$. Then we have
    $$
    {\rm gcd}(f,D_q f){\rm gcd}(g,D_q g)=[z-z_1]^{m+n}_qF(z)G(z).
    $$
It follows from \eqref{fDf.eq}, \eqref{gDg.eq} and Corollary \ref{n-1.cor} that
    $$
    f(z)=[z-z_1]^{m+1}_qF_1(z),
    $$
where $F_1$ is an entire function, and
    $$
    g(z)=[z-z_2]^{n+1}_qG_1(z),
    $$
where $G_1$ is an entire function. Therefore, we have
    $$
    f(z)g(z)=[z-z_1]^{m+n+1}_q(z-q^mz_1)F_1(z)G_1(z).
    $$
It shows that $f$ and $g$ have a common $q$-divisor $z-z_1$, which is a contradiction
with our assumption. We thus proved the first assertion.

Suppose that ${\rm gcd}(f,D_q f)$ and ${\rm gcd}(g,D_q g)$ are not relatively prime.
We set $z-z_0$ is the common divisor of ${\rm gcd}(f,D_q f)$ and ${\rm gcd}(g,D_q g)$. Hence,
there exist entire functions $F_2, G_2$ such that
    $$
    {\rm gcd}(f,D_q f)=(z-z_0)F_2(z),
    $$
and
    $$
    {\rm gcd}(g,D_q g)=(z-z_0)G_2(z).
    $$
It yields by Corollary \ref{n-1.cor} that $f=[z-z_0]^2_qF_3$ and $g=[z-z_0]_q^2G_3$, where $F_3$ and $G_3$ are entire functions. It implies that $f$ and $g$ are not relatively shifting prime, which
is a contradiction to our assumption. We proved our Lemma \ref{q-prime}.
\end{proof}

\end{document}
